\documentclass[11pt]{article}
\usepackage[margin=1in]{geometry}
\usepackage{amsmath,amsthm,amssymb,booktabs}
\usepackage[hidelinks]{hyperref}
\newtheorem{theorem}{Theorem}
\newtheorem{lemma}{Lemma}
\newtheorem{corollary}{Corollary}
\newcommand{\E}{\mathbb E}
\newcommand{\CM}{\mathrm{CM}}
\title{Which independent refresh chains have completely monotone Shannon entropy deficit?}
\author{Leo Gao\\{\small AI-assisted research}}
\date{}
\begin{document}
\maketitle
\begin{abstract}
Consider a finite product space in which coordinate $i$ is independently
redrawn from its stationary law at rate $r_i>0$. Relative entropy to
stationarity decreases, but need not have alternating derivatives of all
orders. We classify the chains for which it does so for every starting law.
After deterministic factors are deleted, the only nontrivial possibilities
are one fair bit, or two fair bits refreshed at equal rates.
The positive two-bit theorem uses a coefficient bound for a pair of
quadratic polynomials. Even an arbitrarily small inequality of the two rates
destroys the universal property. A three-bit example has a certified
negative signed nineteenth derivative. This is a finite-channel result,
not a resolution of the Gaussian heat-flow conjecture.
\end{abstract}
\section{Problem and result}
On $\Omega=\prod_i\Omega_i$ let $\pi=\bigotimes_i\pi_i$, with finite
alphabets and full support. Independently in coordinate $i$, a rate-$r_i$
Poisson clock replaces the current symbol by a fresh $\pi_i$ draw. Thus the
one-coordinate semigroup at time $t$ is
$e^{-r_it}I+(1-e^{-r_it})\Pi_i$, where $\Pi_i$ averages over the coordinate.
For an arbitrary starting law $\mu$ write
\[
 D_\mu(t)=\sum_{x\in\Omega}\mu_t(x)\log\frac{\mu_t(x)}{\pi(x)}.
\]
With natural logarithms, $D_\mu(t)$ is nonnegative and nonincreasing.
Complete monotonicity means $(-1)^kD_\mu^{(k)}(t)\ge0$ for all integers
$k\ge0$ and $t>0$.
\begin{theorem}\label{class}
Every initial law has completely monotone entropy deficit if and only if,
after deletion of one-point coordinates, the space is empty, consists of
one fair bit, or consists of two fair bits with $r_1=r_2$.
\end{theorem}
For uniform stationary bits, the channel is precisely independent binary
symmetric noise with coordinate crossover probabilities $(1-e^{-r_it})/2$.
Entropy under this channel has an extensive literature, including
Samorodnitsky~\cite{sam}; usual comparisons at a fixed noise level do not
settle signs of all time derivatives. Complete monotonicity for several
linear and quadratic observables of reversible chains follows from spectral
positivity~\cite{af}. The nonlinear continuous heat-flow question of
Cheng and Geng~\cite{cg} was disproved by Gu and Sellke~\cite{gs}; the
continuous theorem neither states nor implies the discrete classification
above. We make no claim of identifying the first failing derivative.

\section{A coefficient test}
\begin{lemma}\label{spec}
Suppose $D(t)=\sum_{\lambda}d_\lambda e^{-\lambda t}$ at all sufficiently
large $t$, with distinct nonnegative frequencies, and absolute convergence
at $T$. If $d_{\lambda_0}<0$ for some $\lambda_0>0$, then $D$ is not
completely monotone. Conversely, if all coefficients are nonnegative and
the expansion converges for every $t>0$, it is completely monotone.
\end{lemma}
\begin{proof}
Differentiate termwise for $t>T$ by exponential domination. Evaluate the
signed $k$th derivative at $T+k/\lambda_0$ and multiply by
$(e/\lambda_0)^k$. The result is
\[
 \sum_\lambda d_\lambda e^{-T\lambda}
 \left[(\lambda/\lambda_0)e^{1-\lambda/\lambda_0}\right]^k.
\]
The bracket is in $[0,1]$, equal to one only at $\lambda_0$.
Dominated convergence gives $d_{\lambda_0}e^{-T\lambda_0}<0$.
The converse follows directly by termwise differentiation.
\end{proof}

\section{Why two fair bits work at equal rates}
Normalize the common rate to one and put $z=e^{-t}$. Write the initial
density relative to uniform as $1+ax+by+cxy$, $x,y\in\{-1,1\}$.
The evolved density is $1+azx+bzy+cz^2xy$. Let $\phi(w)=w\log w$.
Its uniform spatial average after this substitution is $D(t)$; the average
is even in $z$. Flipping one input sign if necessary, suppose $c\ge0$.
For $c=0$ the nonnegative coefficients of even degree are immediate
from $\phi(1+u)=u+\sum_{j\ge2}(-1)^ju^j/[j(j-1)]$ and the symmetry
of $ax+by$. Now take $c>0$ and put $w=\sqrt c\,z$. Pairing opposite signs
reduces each even coefficient to half the sum of the corresponding
coefficients of
\[
 \phi(1-Aw+w^2),\quad \phi(1-Bw-w^2),
 \qquad A=(a+b)/\sqrt c,\quad B=(a-b)/\sqrt c.
\]
For even $n\ge4$, the coefficient of $w^n$ in
$\phi((1-rw)(1-sw))$, multiplied by $n(n-1)(n-2)/2$, is
\begin{equation}\label{roots}
 \tfrac12\big[(n-2)(r^n+s^n)-nrs(r^{n-2}+s^{n-2})\big].
\end{equation}
This follows by multiplying the two Taylor series for
$\log(1-rw)+\log(1-sw)$ by the quadratic.
For the minus quadratic $rs=-1$, the roots in this factorization have
absolute values $e^y,e^{-y}$, so \eqref{roots} is
$(n-2)\cosh(ny)+n\cosh((n-2)y)\ge2n-2$.
For the plus quadratic $rs=1$. If its factors are complex conjugates,
\eqref{roots} is $(n-2)\cos(n\theta)-n\cos((n-2)\theta)\ge-(2n-2)$.
If they are real, it is $(n-2)\cosh(ny)-n\cosh((n-2)y)\ge-2$:
this is even in $y$ and increasing for $y\ge0$, by differentiation.
Thus the paired even coefficients are nonnegative. Degree two in $z$ is
$(a^2+b^2)/2$, and the constant is zero.

The series indeed converges for every $|z|<1$: the original four
quadratics in $z$ are strictly positive for real $-1<z<1$ because the
channel with crossover $(1-z)/2$ has positive entries.
Nonreal zeros can occur only in the positive-quadratic case, and have
modulus $1/\sqrt c\ge1$, since $|c|=|\E_\mu XY|\le1$.
This gives analytic logarithms on the open disk. Apply Lemma~\ref{spec}.
A single fair bit is a marginal case obtained by adjoining an independent
stationary fair bit; the empty product has identically zero deficit.

\section{Obstructions}
The class of counterexamples is closed under adjoining independent
stationary coordinates. First suppose one site admits an event of
probability $p\ne1/2$, $0<p<1$. Then $h=\mathbf1_A-p$ has zero mean
and third moment $p(1-p)(1-2p)\ne0$. For the single-site starting density
$1+\epsilon h$, choose small $\epsilon$ so that $\epsilon^3\E h^3>0$.
The coefficient at frequency $3r$ is
$-\epsilon^3\E h^3/6<0$. This applies to all nontrivial sites other than
fair binary sites, since at least one singleton then has probability
different from $1/2$.

For three fair bits of equal rate one, put
$g=xy+xz+yz$ (here $x,y,z$ are signs). It is $3$ with probability
$1/4$, $-1$ otherwise. The semigroup sends $g$ to $e^{-2t}g$.
For $0<\epsilon<1$, $1+\epsilon g$ is a positive density and
the frequency-six coefficient is $-\epsilon^3\E g^3/6=-\epsilon^3$.
This gives failures arbitrarily close to equilibrium. The verification
script also certifies
\[
 -.095<2^{-19}(-1)^{19}D^{(19)}(\log 16)<-.094
 \quad(\epsilon=1/2).
\]

It remains to analyze two unequal bit rates $r<s$. Use
$X=e^{-rt},Y=e^{-st}$ and the starting law $1+ax+by+vabxy$.
Expand the resulting $D$ in powers of $(aX,bY)$ near the origin.
Independent sign symmetry kills odd degrees in each variable.
For $j\ge2$ even, the coefficient of $(bY)^j$ is
\[
 \frac1{j(j-1)}\E_x\frac{(1+vaXx)^j}{(1+aXx)^{j-1}}.
\]
Indeed condition on $x$, expand $\phi$ about $1+aXx$, and average over $y$.
Consequently for even $i$ the coefficient $d_{ij}$ of $(aX)^i(bY)^j$
is $[w^i](1+vw)^j(1+w)^{1-j}/[j(j-1)]$.
For $j=0$ and $i\ge2$ even it is $1/[i(i-1)]$.
In particular,
\[
 d_{44}=\frac{(v-1)^2(v^2-10v+15)}{12}=-45/64
 \quad\text{at }v=5/2.
\]
Fix this $v$. If $s/r$ is irrational, the coefficient
$d_{44}a^4b^4$ is alone at frequency $4r+4s$.
If $r:s=p:q$ in lowest terms with $q\ge3$, collisions must have
$(i,j)=(4+2qk,4-2pk)$ with integer $k$. Nonnegative $i$ forces $k\ge0$.
For fixed $0<b<1/4$ and sufficiently small $a>0$ the frequency therefore
has coefficient $d_{44}a^4b^4+O(a^{4+2q})<0$.
All four initial density values are positive for sufficiently small $a$.
These power series converge absolutely at sufficiently large $t$,
so aggregating coincident frequencies is legitimate.

The remaining unequal ratio is $1:2$. In units where $r=1$,
put $a=1/8$, $b=(3/2)a^2=:Ka^2$; since $a+b+vab<1$, the density
is positive. The frequency-twelve coefficient divided by $a^{12}$ is
\[
 d_{12,0}+d_{8,2}K^2+d_{4,4}K^4+d_{0,6}K^6
 =\frac1{132}+\frac98K^2-\frac{45}{64}K^4+\frac1{30}K^6
 =-\frac{108313}{168960}.
\]
Lemma~\ref{spec} finishes the proof of Theorem~\ref{class}.
\section{Scope and verification}
The theorem concerns independent full redraws in continuous time and
universal quantification over starting laws. It does not classify
general reversible chains, Boolean functions maximizing mutual information,
or the Gaussian heat equation. For independent inputs the deficit adds
across coordinates, so higher-dimensional failure for fair bits genuinely
uses dependence. The checker \texttt{verify.py} uses exact rational
arithmetic to verify the $1:2$ witness by two independent expansions,
test low-degree coefficients on a finite two-bit grid, and bound the
nineteenth-derivative witness with a geometric-series tail. These checks
supplement, not replace, the proof. A targeted literature search found
no matching full classification; that is not a priority determination.
\begin{thebibliography}{9}
\bibitem{af} D. Aldous and J. A. Fill, \emph{Reversible Markov Chains
and Random Walks on Graphs}, Ch.~3, \S3.5,
\href{https://www.stat.berkeley.edu/~aldous/RWG/Book_Ralph/Ch3.S5.html}{Complete monotonicity}.
\bibitem{cg} F. Cheng and Y. Geng, Higher order derivatives in Costa's
entropy power inequality, \emph{IEEE Transactions on Information Theory}
61 (2015), 5892--5905. \href{https://arxiv.org/abs/1409.5543}{arXiv:1409.5543}.
\bibitem{gs} Y. Gu and M. Sellke, A counterexample to the Gaussian
completely monotone conjecture (2026),
\href{https://arxiv.org/abs/2605.11656}{arXiv:2605.11656}.
\bibitem{sam} A. Samorodnitsky, On the entropy of a noisy function (2015),
\href{https://arxiv.org/abs/1508.01464}{arXiv:1508.01464}.
\end{thebibliography}
\end{document}
